\documentclass[11pt]{article}

\usepackage[a4paper,margin=27mm]{geometry}
\usepackage{amsmath,amssymb,amsthm,mathtools}
\usepackage[hidelinks]{hyperref}
\usepackage{microtype}
\usepackage{enumitem}

\newtheorem{definition}{Definition}[section]
\newtheorem{proposition}[definition]{Proposition}
\newtheorem{theorem}[definition]{Theorem}
\newtheorem{corollary}[definition]{Corollary}
\newtheorem{example}[definition]{Example}
\newtheorem{remark}[definition]{Remark}

\newcommand{\X}{\mathcal X}
\newcommand{\Acal}{\mathcal A}
\newcommand{\Bcal}{\mathcal B}

\hypersetup{
  pdftitle={WR-III: Geometry of Admissible World Fibres},
  pdfauthor={A. A. Malachevsky},
  pdfsubject={World Realizability and Epistemic Horizons (WREH), WR-III},
  pdfkeywords={world fibres, inverse images, identifiability, rigidity, refinement, set-valued maps, realizability walls}
}

\title{\textbf{WR-III Preprint Review v0.5}\\[2pt]
\normalsize Review copy -- for audit only; not for citation\\[6pt]
\large Geometry of Admissible World Fibres:\\
Refinement, Rigidity and Realizability Walls}
\author{A. A. Malachevsky\\
ORCID: 0009-0008-6009-3196\\
World Realizability \& Epistemic Horizons (WREH)\\
Research Programme \& Community}
\date{8 October 2026}

\begin{document}
\maketitle

\begin{abstract}
WR-III studies the geometry of nonempty inverse fibres left after the
identifiability and realizability questions of WR-I and WR-II. For a declared
geometric realization carrier and a response map indexed by a protocol bundle,
we introduce the inverse-fibre correspondence over the realizable response
image. We prove monotone fibre shrinkage under protocol refinement, establish
upper hemicontinuity on compact carriers, characterize lower hemicontinuity by
openness of the response map, and derive upper semicontinuity of fibre diameter
on compact metric carriers. These results yield open finite-resolution
rigidity loci. In the smooth proper setting, Ehresmann's theorem supplies local
triviality over regular values, motivating an intrinsic structural wall as the
failure locus of local fibre triviality. This wall is kept distinct from the
boundary of the attainable response image. The framework is model-relative:
no response-space wall is identified with a physical boundary or phase
transition without an additional bridge. We further prove a refinement
asymmetry: exact and finite-resolution rigidity are monotone under added
protocols, whereas structural bifurcation walls need not be. A planar
range-only localization benchmark realizes this contrast explicitly through
the one-anchor, two-anchor, and three-noncollinear-anchor ladder.
\end{abstract}

\section{Claim ceiling and scope}

WR-III is a geometric continuation of WR-I and WR-II. Its claims are
model-relative and conditional on the declared realization carrier, response
maps, topology, metric, and smooth or stratified structure.

WR-III does not claim that:
\begin{itemize}
\item a response fibre is an ontological ensemble of actually existing worlds;
\item fibre diameter is physically meaningful without a declared metric on the realization carrier;
\item protocol refinement is physical time, dynamics, causal evolution, or RG flow;
\item a bifurcation value of a response map is a thermodynamic or quantum phase transition;
\item a WREH ``wall'' is a wall of spacetime or of the Universe;
\item the open/closed-map criteria, Ehresmann triviality, relative polar geometry, Hardt triviality, or Reeb-space constructions are new;
\item the classical trilateration formulas in the localization benchmark are new.
\end{itemize}

The paper's narrower target is the organization of nonempty response fibres
downstream of WR-I identifiability and WR-II realizability, with special
attention to protocol-refinement shrinkage, metric rigidity, conditional rank,
and bifurcation geometry.
\section{Upstream interface}

WR-III begins after WR-I and WR-II. Let \(\X\) be the declared geometric
realization carrier and let
\[
R_{\Acal}:\X\to Y_{\Acal}
\]
be the response map associated with an admissible protocol bundle \(\Acal\).
Its realizable response image is
\[
E_{\Acal}:=R_{\Acal}(\X).
\]

The carrier \(\X\) may be the original admissible world class or a justified
gauge-reduced carrier. WR-III does not assume that a gauge quotient exists, is
Hausdorff, metric, or smooth unless those properties are stated.

\begin{definition}[Admissible world fibre]
For \(e\in E_{\Acal}\), define
\[
\Phi_{\Acal}(e):=R_{\Acal}^{-1}(e).
\]
\end{definition}

Every fibre considered in WR-III is nonempty by construction. Whether a
response is realizable belongs upstream to WR-II.

\section{Refinement geometry}

Suppose \(\Acal\subseteq\Bcal\) and
\[
R_{\Acal}
=
\rho_{\Bcal\Acal}\circ R_{\Bcal}.
\]

\begin{proposition}[Refinement inclusion]
For every \(e_{\Bcal}\in E_{\Bcal}\),
\[
\Phi_{\Bcal}(e_{\Bcal})
\subseteq
\Phi_{\Acal}\bigl(\rho_{\Bcal\Acal}(e_{\Bcal})\bigr).
\]
\end{proposition}

\begin{proof}
If \(x\in\Phi_{\Bcal}(e_{\Bcal})\), then
\[
R_{\Bcal}(x)=e_{\Bcal}.
\]
Applying \(\rho_{\Bcal\Acal}\) gives
\[
R_{\Acal}(x)=\rho_{\Bcal\Acal}(e_{\Bcal}),
\]
so \(x\) lies in the coarser fibre.
\end{proof}

Assume now that \((\X,d)\) is metric.

\begin{definition}[Fibre diameter]
For nonempty \(F\subseteq\X\), define
\[
\operatorname{diam}F
=
\sup\{d(x,x'):x,x'\in F\}.
\]
Set
\[
D_{\Acal}(e)
:=
\operatorname{diam}\Phi_{\Acal}(e).
\]
\end{definition}

\begin{corollary}[Diameter monotonicity under refinement]
For \(\Acal\subseteq\Bcal\),
\[
D_{\Bcal}(e_{\Bcal})
\le
D_{\Acal}\bigl(\rho_{\Bcal\Acal}(e_{\Bcal})\bigr).
\]
\end{corollary}

\begin{remark}
This monotonicity is informational. It does not imply a physical time arrow,
dynamical contraction, entropy law, or renormalization flow.
\end{remark}

\section{Rigidity under refinement and conditional rank}

Let \(\Acal\subseteq\Bcal\) with response restriction
\[
R_{\Acal}=\rho_{\Bcal\Acal}\circ R_{\Bcal}.
\]

\begin{definition}[Rigidity loci]
For \(\varepsilon>0\), define
\[
\mathcal R_{\Acal}^{\varepsilon}
=
\{e\in E_{\Acal}:D_{\Acal}(e)<\varepsilon\},
\]
and let \(\mathcal R_{\Acal}^{0}\) denote the exact-rigidity locus.
\end{definition}

\begin{theorem}[Pulled-back rigidity monotonicity]
\label{thm:rigidity-pullback}
For every \(\varepsilon\ge0\),
\[
E_{\Bcal}\cap
\rho_{\Bcal\Acal}^{-1}
(\mathcal R_{\Acal}^{\varepsilon})
\subseteq
\mathcal R_{\Bcal}^{\varepsilon}.
\]
Thus protocol refinement cannot destroy exact or finite-resolution rigidity
already present at the coarser response.
\end{theorem}

\begin{proof}
For \(\varepsilon>0\), this is immediate from fibre inclusion and diameter
monotonicity. For \(\varepsilon=0\), a nonempty refined fibre contained in
a singleton coarse fibre is the same singleton.
\end{proof}

\begin{definition}[Rigidity gain]
For a metric realization carrier define
\[
G_{\Bcal\mid\Acal}(e_{\Bcal})
:=
D_{\Acal}(\rho_{\Bcal\Acal}(e_{\Bcal}))
-D_{\Bcal}(e_{\Bcal})
\ge0.
\]
This is a metric- and protocol-dependent quantity measuring how much the
added response block reduces fibre diameter.
\end{definition}

The smooth finite-dimensional case admits a sharper local criterion. Suppose
\[
R_{\Bcal}=(R_{\Acal},S):\X\to Y_{\Acal}\times Z,
\]
where \(S:\X\to Z\) is the added response block.

\begin{theorem}[Conditional-rank formula for an added response block]
\label{thm:conditional-rank}
Let \(x\in\X\) be a point at which
\[
dR_{\Acal,x}:T_x\X\to T_{R_{\Acal}(x)}Y_{\Acal}
\]
is surjective. Then
\[
\operatorname{rank}dR_{\Bcal,x}
=
\dim Y_{\Acal}
+
\operatorname{rank}
\left(
dS_x\big|_{\ker dR_{\Acal,x}}
\right).
\]
In particular, \(R_{\Bcal}\) is a submersion at \(x\) if and only if
\[
dS_x\big|_{\ker dR_{\Acal,x}}:
\ker dR_{\Acal,x}\to T_{S(x)}Z
\]
is surjective.
\end{theorem}

\begin{proof}
Since \(dR_{\Acal,x}\) is surjective, choose a linear splitting
\[
T_x\X=H_x\oplus\ker dR_{\Acal,x}
\]
such that \(dR_{\Acal,x}|_{H_x}\) is an isomorphism onto
\(T_{R_{\Acal}(x)}Y_{\Acal}\). The image of
\[
dR_{\Bcal,x}=(dR_{\Acal,x},dS_x)
\]
contains a graph over the entire coarse-response tangent space, while the
additional independent directions in the \(Z\)-factor are exactly
\[
dS_x(\ker dR_{\Acal,x}).
\]
Hence the rank is the coarse target dimension plus the rank of this
conditional derivative. Surjectivity onto the product is therefore
equivalent to surjectivity of the conditional derivative.
\end{proof}

\begin{definition}[Conditional rank-defect set]
On the coarse-regular locus define
\[
\mathcal C_{\Bcal\mid\Acal}
=
\left\{
x:
dR_{\Acal,x}\text{ is surjective and }
\operatorname{rank}
\left(dS_x|_{\ker dR_{\Acal,x}}\right)
<\dim Z
\right\}.
\]
\end{definition}

\begin{corollary}[Container for refinement-created walls]
\label{cor:new-wall-container}
Assume \(R_{\Bcal}=(R_{\Acal},S)\) is smooth and proper. Let
\(e_{\Bcal}=(e_{\Acal},z)\) be a refined response whose coarse component
\(e_{\Acal}\) is a regular value of \(R_{\Acal}\). If
\(e_{\Bcal}\) lies on the structural wall of \(R_{\Bcal}\), then
\[
e_{\Bcal}\in
R_{\Bcal}(\mathcal C_{\Bcal\mid\Acal}).
\]
Thus a new refined wall above a coarse-regular response requires loss of
conditional rank in the added response block along the coarse fibre.
\end{corollary}

\begin{proof}
Proper smooth structural-wall points are critical values of
\(R_{\Bcal}\). Hence some
\(x\in R_{\Bcal}^{-1}(e_{\Bcal})\) is critical for the refined map.
Because \(e_{\Acal}\) is a regular value, \(dR_{\Acal,x}\) is
surjective. Theorem~\ref{thm:conditional-rank} then forces conditional
rank loss at \(x\).
\end{proof}

\begin{remark}[Interpretation firewall]
The conditional-rank criterion is differential geometry of an observation
map. It is not by itself a phase-transition criterion, dynamical instability,
or physical criticality statement.
\end{remark}

\section{Upper hemicontinuity and the closed-map criterion}

A correspondence
\[
\Phi:E\rightrightarrows\X
\]
is upper hemicontinuous at \(e\) if for every open \(U\subseteq\X\) with
\[
\Phi(e)\subseteq U,
\]
there is a neighborhood \(V\) of \(e\) such that
\[
\Phi(e')\subseteq U
\qquad(e'\in V).
\]

\begin{theorem}[Closed-map criterion]
\label{thm:closed-map}
Let \(R:\X\to E\) be a surjection. The inverse correspondence
\[
\Phi(e)=R^{-1}(e)
\]
is upper hemicontinuous if and only if \(R\) is a closed map.
\end{theorem}

\begin{proof}
Assume first that \(R\) is closed. Let \(U\subseteq\X\) be open with
\(\Phi(e)\subseteq U\). Then \(C=\X\setminus U\) is closed and
\(e\notin R(C)\). Since \(R(C)\) is closed,
\[
V:=E\setminus R(C)
\]
is an open neighborhood of \(e\), and every fibre above \(V\) is contained
in \(U\).

Conversely, suppose the inverse correspondence is upper hemicontinuous and let
\(C\subseteq\X\) be closed. If \(e\notin R(C)\), then
\[
\Phi(e)\subseteq\X\setminus C.
\]
Upper hemicontinuity gives a neighborhood \(V\) of \(e\) whose fibres avoid
\(C\). Hence \(V\cap R(C)=\varnothing\), so \(R(C)\) is closed.
\end{proof}

\begin{corollary}[Compact-carrier upper hemicontinuity]
\label{thm:uhc}
Let \(\X\) be compact, let \(E\) be Hausdorff, and let
\[
R:\X\to E
\]
be continuous and surjective. Then the inverse correspondence has nonempty
compact values and is upper hemicontinuous.
\end{corollary}

\begin{proof}
A continuous map from a compact space to a Hausdorff space is closed. Apply
Theorem~\ref{thm:closed-map}. Each fibre is the preimage of a closed singleton
and is therefore compact.
\end{proof}
\section{Lower hemicontinuity and openness}

A correspondence \(\Phi:E\rightrightarrows\X\) is lower hemicontinuous if for
every open \(U\subseteq\X\), the set
\[
\{e\in E:\Phi(e)\cap U\neq\varnothing\}
\]
is open.

\begin{theorem}[Openness criterion]
\label{thm:open-map}
Let \(R:\X\to E\) be a surjection. Its inverse correspondence
\[
\Phi(e)=R^{-1}(e)
\]
is lower hemicontinuous if and only if \(R\) is an open map.
\end{theorem}

\begin{proof}
For every open \(U\subseteq\X\),
\[
R(U)
=
\{e\in E:\Phi(e)\cap U\neq\varnothing\}.
\]
Thus the defining condition for lower hemicontinuity is exactly openness of
\(R\).
\end{proof}

\begin{corollary}
If \(\X\) is compact, \(E\) is Hausdorff, and \(R\) is continuous, surjective,
and open, then its inverse correspondence is both upper and lower
hemicontinuous.
\end{corollary}

\section{Fibre-diameter stability}

\begin{theorem}[Diameter upper semicontinuity]
\label{thm:diam-usc}
Let \((\X,d)\) be compact metric, let \(E\) be metric, and let
\[
R:\X\to E
\]
be continuous and surjective. Then
\[
D(e):=\operatorname{diam}R^{-1}(e)
\]
is upper semicontinuous.
\end{theorem}

\begin{proof}
Let \(e_n\to e\). Choose
\[
x_n,x_n'\in R^{-1}(e_n)
\]
with
\[
d(x_n,x_n')
\ge D(e_n)-\frac1n.
\]
Pass to a subsequence realizing the limsup. Compactness gives
\[
x_n\to x,\qquad x_n'\to x'.
\]
Continuity gives \(R(x)=R(x')=e\), and therefore
\[
\limsup_{n\to\infty}D(e_n)
\le d(x,x')
\le D(e).
\]
\end{proof}

\begin{corollary}[Open finite-resolution rigidity loci]
For every \(\varepsilon>0\),
\[
\mathcal R_\varepsilon
=
\{e\in E:D(e)<\varepsilon\}
\]
is open.
\end{corollary}

\begin{corollary}[Exact rigidity is a \(G_\delta\) locus]
The exact rigidity locus
\[
\mathcal R_0
=
\{e\in E:D(e)=0\}
\]
satisfies
\[
\mathcal R_0
=
\bigcap_{n=1}^{\infty}\mathcal R_{1/n}.
\]
\end{corollary}

\begin{theorem}[Hausdorff continuity under an open response map]
\label{thm:hausdorff-cont}
Let \((\X,d)\) be compact metric, let \(E\) be metric, and let
\[
R:\X\to E
\]
be continuous, surjective, and open. Then the compact fibres
\[
e\longmapsto R^{-1}(e)
\]
vary continuously in the Hausdorff metric. Consequently, \(D(e)\) is
continuous.
\end{theorem}

\begin{proof}
By Theorems~\ref{thm:closed-map} and \ref{thm:open-map}, the fibre
correspondence is upper and lower hemicontinuous. Let \(e_n\to e\).
Upper hemicontinuity gives, for every \(\varepsilon>0\),
\[
R^{-1}(e_n)\subseteq N_\varepsilon(R^{-1}(e))
\]
for all sufficiently large \(n\). Conversely, cover the compact fibre
\(R^{-1}(e)\) by finitely many balls \(B(x_j,\varepsilon/2)\). Lower
hemicontinuity implies that for all large \(n\), every such ball meets
\(R^{-1}(e_n)\). Hence
\[
d_H(R^{-1}(e_n),R^{-1}(e))\to0.
\]
For nonempty compact \(A,B\subseteq\X\),
\[
|\operatorname{diam}A-\operatorname{diam}B|\le2d_H(A,B),
\]
so \(D(e_n)\to D(e)\).
\end{proof}
\begin{remark}
Exact rigidity need not be open. Unique identification can occur at a
structural singularity.
\end{remark}

\section{Structural regularity, bifurcation values, and walls}

The word wall is used here in a strictly mathematical, model-relative sense.

\begin{definition}[Structural regular response]
A response \(e\in E\) is structurally regular if there exists an open
neighborhood \(U\subseteq E\) of \(e\) such that
\[
R^{-1}(U)\to U
\]
is a locally trivial topological fibre bundle.
\end{definition}

Let \(E_{\mathrm{reg}}^{\mathrm{str}}\) be the union of all structurally regular
neighborhoods.

\begin{definition}[WREH structural wall / bifurcation locus]
The structural wall is
\[
\Sigma_{\mathrm{str}}
=
E\setminus E_{\mathrm{reg}}^{\mathrm{str}}.
\]
\end{definition}

\begin{remark}[Standard terminology]
In smooth and singularity-theory settings, the failure-of-local-triviality
locus is commonly called the bifurcation set or set of atypical values.
WR-III does not claim novelty for that notion; the word ``wall'' is retained
only as WREH programme vocabulary.
\end{remark}

\begin{remark}[No codimension claim]
Neither the structural wall nor the combined realizability wall is assumed to
be a hypersurface, manifold, or codimension-one set.
\end{remark}
\begin{definition}[Attainable-response boundary]
Let \(M\) be a declared response manifold or tame response stratum in which
\(E\subseteq M\) is full-dimensional. Define
\[
\partial_M E
\]
to be the topological boundary of \(E\) in \(M\).
\end{definition}

\begin{definition}[Combined WREH realizability-wall diagnostic]
When such a full-dimensional ambient response space \(M\) has been justified,
define
\[
\Sigma_{\mathrm{real}}
=
\Sigma_{\mathrm{str}}\cup\partial_M E.
\]
\end{definition}

\begin{remark}[Ambient-space firewall]
If \(E\) is lower-dimensional in an arbitrarily chosen ambient product space,
its ambient boundary may equal all of \(E\). In that situation the combined
diagnostic is not formed until an intrinsic or full-dimensional response
manifold/stratum is declared. The internal structural wall
\(\Sigma_{\mathrm{str}}\) remains well defined on \(E\) itself.
\end{remark}

\begin{remark}
The two pieces have different meanings. The boundary separates attainable
from unattainable responses in the declared full-dimensional response model.
The internal structural wall records failure of local fibre triviality among
attainable responses. Neither is automatically a physical boundary.
\end{remark}

\section{Smooth proper response maps}

Assume \(\X\) and \(Y\) are smooth manifolds without boundary and
\[
R:\X\to Y
\]
is smooth and proper. Let \(E=R(\X)\).

\begin{theorem}[Regular values are structurally regular]
If \(e\in E\) is a regular value of \(R\), then \(e\) is structurally regular.
Consequently,
\[
\Sigma_{\mathrm{str}}
\subseteq
\operatorname{CritVal}(R)\cap E.
\]
\end{theorem}

\begin{proof}
The critical set is closed. Since a proper map between manifolds is closed,
its critical-value set is closed.

Choose \(x\in R^{-1}(e)\). Because \(e\) is a regular value, the
differential \(dR_x\) is surjective. The submersion theorem therefore gives
an open neighborhood \(U_0\subseteq Y\) of \(e\) contained in the image
\(E\). Since the critical-value set is closed and does not contain \(e\),
there is an open neighborhood \(U_1\subseteq Y\) of \(e\) containing no
critical values. Set \(U=U_0\cap U_1\). Then
\[
R:R^{-1}(U)\to U
\]
is surjective, proper, and a submersion. Ehresmann's fibration theorem implies
local triviality.
\end{proof}

\begin{corollary}[Constant smooth fibre type off the wall]
On each connected component of the regular-value region, all fibres are
diffeomorphic.
\end{corollary}

\begin{remark}
Critical values only contain the possible structural wall under these
hypotheses. A critical value need not change every fibre invariant of interest.
\end{remark}

\section{Minimal example}

\begin{example}[Height response on the sphere]
Let
\[
\X=S^2\subset\mathbb R^3,
\qquad
R(x,y,z)=z.
\]
Then
\[
E=[-1,1].
\]
For \(-1<e<1\), the fibre is a circle. At \(e=\pm1\), the fibre is a
singleton.

With the ambient Euclidean metric,
\[
D(e)=2\sqrt{1-e^2}.
\]
Thus exact rigidity occurs precisely at the two structural wall points.
\end{example}

\begin{remark}[Rigidity is not regularity]
The example shows that exact rigidity and structural regularity are logically
distinct. A singleton fibre can occur where the local fibre type becomes
singular.
\end{remark}

\section{Refinement can create or resolve structural walls}

The rigidity loci of Theorem~\ref{thm:rigidity-pullback} are monotone under
refinement after pullback. Structural walls behave differently.

\begin{example}[Refinement creates an internal cusp wall]
Let
\[
\X=\mathbb R^2,
\qquad
R_{\Acal}(x,y)=x.
\]
The coarse map is a globally trivial bundle with no structural wall.
Add the response
\[
S(x,y)=y^3-xy
\]
and define
\[
R_{\Bcal}(x,y)=(x,y^3-xy).
\]
The restriction map is projection onto the first coordinate, so this is a
genuine refinement.

The refined map is proper and surjective onto \(\mathbb R^2\). Its Jacobian
determinant is
\[
3y^2-x.
\]
Therefore the conditional rank-defect set is
\[
x=3y^2.
\]
Its image is the semicubical cusp
\[
27v^2=4u^3,
\qquad u\ge0.
\]
Across this cusp the cubic equation
\[
y^3-uy=v
\]
changes from one real solution to three real solutions. Hence the cusp is the
refined structural bifurcation wall, even though the coarse response map had
no wall at all.
\end{example}

\begin{example}[Refinement resolves a coarse branch wall]
Let
\[
\X=[-1,1]\times S^1
\]
and define the coarse response
\[
R_{\Acal}(x,\theta)=x^2,
\qquad
E_{\Acal}=[0,1].
\]
For \(e>0\), the fibre is the disjoint union of two circles, while at
\(e=0\) it is one circle. Thus \(0\) is a structural wall.

Now refine by recording the branch label \(x\):
\[
R_{\Bcal}(x,\theta)=(x^2,x).
\]
Its response image is the parabola arc
\[
E_{\Bcal}
=
\{(u,v):u=v^2,\ |v|\le1\}.
\]
Projection \((u,v)\mapsto u\) recovers the coarse response. But
\(R_{\Bcal}:\X\to E_{\Bcal}\) is globally a trivial bundle with fibre
\(S^1\), via the homeomorphism
\[
(x,\theta)\longmapsto ((x^2,x),\theta).
\]
Hence the refined structural wall is empty.
\end{example}

\begin{theorem}[Refinement asymmetry]
Structural walls are not monotone under protocol refinement: a refinement can
create new structural walls or resolve coarse structural walls. In contrast,
exact and finite-resolution rigidity loci obey the monotonicity relation of
Theorem~\ref{thm:rigidity-pullback}.
\end{theorem}

\begin{proof}
The two preceding examples establish creation and resolution of walls.
Rigidity monotonicity was proved in Theorem~\ref{thm:rigidity-pullback}.
\end{proof}

\begin{remark}
This asymmetry is a programme-level organizing contrast in WR-III: ambiguity
shrinks monotonically under added information, while the singularity structure
of the response map need not. No theorem-priority claim is made for the
underlying singularity mechanisms.
\end{remark}

\section{Application benchmark: planar range-only localization}
\label{sec:range-localization}

Range-based localization is a standard inverse problem in sensor networks and
multilateration: an unknown point is reconstructed from distances to anchors of
known position \cite{AspnesEtAl2006,AndersonEtAl2010}. We use its simplest
planar form as a WR-III benchmark because protocol refinement changes both
fibre size and bifurcation geometry in closed form.

Let
\[
\X=\mathbb R^2,
\qquad
p=(x,y),
\]
with the Euclidean metric. Fix two anchors
\[
a_1=(-a,0),
\qquad
a_2=(a,0),
\qquad
a>0,
\]
and define squared-range protocols
\[
r_i(p)=\|p-a_i\|^2.
\]

\subsection{One anchor: circular ambiguity}

For the coarse protocol family
\[
\Acal=\{1\},
\qquad
R_{\Acal}=r_1,
\]
the realizable image is
\[
E_{\Acal}=[0,\infty).
\]
For \(u>0\),
\[
\Phi_{\Acal}(u)
=
\{p:\|p-a_1\|^2=u\}
\]
is a circle of radius \(\sqrt u\), while
\[
\Phi_{\Acal}(0)=\{a_1\}.
\]
Hence
\[
D_{\Acal}(u)=2\sqrt u.
\]

Over \((0,\infty)\) the fibres form the trivial circle bundle
\[
(0,\infty)\times S^1.
\]
At \(u=0\) the circle collapses to a point, so
\[
\Sigma_{\Acal}^{\mathrm{str}}=\{0\}.
\]

\subsection{Two anchors: reflection ambiguity and a fold wall}

Add the second squared-range protocol:
\[
\Bcal=\{1,2\},
\qquad
R_{\Bcal}(p)
=
(r_1(p),r_2(p))
=
(u,v).
\]
From
\[
u=(x+a)^2+y^2,
\qquad
v=(x-a)^2+y^2,
\]
we obtain
\[
x=\frac{u-v}{4a}
\]
and
\[
y^2
=
q(u,v)
:=
\frac{u+v}{2}
-a^2
-\frac{(u-v)^2}{16a^2}.
\]

Therefore the realizable response image is
\[
E_{\Bcal}
=
\left\{
(u,v)\in[0,\infty)^2:
q(u,v)\ge0
\right\}.
\]

\begin{proposition}[Exact two-anchor fibre geometry]
For \((u,v)\in E_{\Bcal}\):
\[
\Phi_{\Bcal}(u,v)
=
\begin{cases}
\left\{
\left(\dfrac{u-v}{4a},\pm\sqrt{q(u,v)}\right)
\right\},
&
q(u,v)>0,
\\[10pt]
\left\{
\left(\dfrac{u-v}{4a},0\right)
\right\},
&
q(u,v)=0.
\end{cases}
\]
Consequently,
\[
D_{\Bcal}(u,v)
=
2\sqrt{q(u,v)}.
\]
\end{proposition}

\begin{proof}
The difference \(u-v\) determines \(x\), and the sum then determines \(y^2\).
The two signs of \(y\) give the mirror ambiguity across the anchor baseline.
The two-point separation is \(2|y|\).
\end{proof}

Thus the refinement from one to two anchors reduces a one-dimensional circular
fibre to at most two points.

For a realized response generated by \(p=(x,y)\),
\[
D_{\Acal}(r_1(p))
=
2\sqrt{(x+a)^2+y^2},
\]
while
\[
D_{\Bcal}(R_{\Bcal}(p))
=
2|y|.
\]
The rigidity gain is therefore
\[
G_{\Bcal\mid\Acal}(R_{\Bcal}(p))
=
2\left(
\sqrt{(x+a)^2+y^2}-|y|
\right)
\ge0.
\]

\begin{theorem}[Two-anchor bifurcation wall]
The structural wall of the two-anchor response map is
\[
\Sigma_{\Bcal}^{\mathrm{str}}
=
\{(u,v)\in E_{\Bcal}:q(u,v)=0\}.
\]
For \(q>0\), \(R_{\Bcal}\) is a two-sheeted covering, one sheet with \(y>0\)
and one with \(y<0\). On \(q=0\), the two sheets merge and the fibre becomes
a singleton.
\end{theorem}

\begin{proof}
On \(q>0\), the explicit inverse formulas define two smooth branches
\[
(u,v)\longmapsto
\left(
\frac{u-v}{4a},
\pm\sqrt{q(u,v)}
\right),
\]
so the restriction is locally trivial with two-point fibre.

On \(q=0\), nearby interior fibres contain two points whereas the wall fibre
contains one point, so local fibre triviality fails.
\end{proof}

The Jacobian of the refined response is
\[
dR_{\Bcal}
=
\begin{pmatrix}
2(x+a) & 2y\\
2(x-a) & 2y
\end{pmatrix},
\]
with determinant
\[
\det dR_{\Bcal}=8ay.
\]
Hence the critical set is precisely the anchor baseline
\[
y=0,
\]
and its image is the fold wall \(q=0\).

The conditional-rank criterion of
Theorem~\ref{thm:conditional-rank} recovers the same geometry directly.
Away from the first anchor \(a_1\), the coarse protocol \(r_1\) is regular and
a tangent vector to its circular fibre is
\[
\tau=(-y,x+a).
\]
For the added protocol \(S=r_2\),
\[
dr_2(\tau)=4ay.
\]
Thus the added range loses conditional rank exactly on
\[
y=0.
\]
The point \(a_1\), where the coarse protocol itself is singular, contributes
the inherited endpoint
\[
R_{\Bcal}(a_1)=(0,4a^2).
\]
Together these give the full fold wall.

\subsection{Three noncollinear anchors: exact localization}

Now add a third anchor
\[
a_3=(0,b),
\qquad
b\neq0,
\]
with squared-range response
\[
r_3(p)=\|p-a_3\|^2.
\]
Let
\[
R_{\mathcal C}(p)
=
(r_1(p),r_2(p),r_3(p))
=
(u,v,w).
\]

\begin{theorem}[Three-anchor exact rigidity]
For three noncollinear anchors as above,
\[
R_{\mathcal C}:\mathbb R^2\to E_{\mathcal C}
\]
is an embedding. Every response fibre is a singleton:
\[
D_{\mathcal C}\equiv0.
\]
Consequently the structural wall relative to the intrinsic response image
\(E_{\mathcal C}\) is empty.
\end{theorem}

\begin{proof}
The first two responses determine
\[
x=\frac{u-v}{4a}
\]
and
\[
s:=x^2+y^2
=
\frac{u+v}{2}-a^2.
\]
The third response satisfies
\[
w=s+b^2-2by,
\]
hence
\[
y=
\frac{s+b^2-w}{2b}.
\]
Thus the response determines \(p=(x,y)\) uniquely.

The map is an immersion. When \(y\neq0\), the first two range gradients already
have rank two. When \(y=0\), the third gradient
\[
\nabla r_3=(2x,-2b)
\]
supplies the missing transverse direction because \(b\neq0\).
Finally, the squared ranges diverge as \(\|p\|\to\infty\), so the map is
proper. A proper injective immersion is an embedding.
\end{proof}

This final refinement removes the two-anchor reflection ambiguity and resolves
the fold wall. The sequence
\[
\{r_1\}
\subset
\{r_1,r_2\}
\subset
\{r_1,r_2,r_3\}
\]
therefore exhibits, in a standard inverse problem,
\[
\text{circle fibre}
\longrightarrow
\text{mirror pair / fold wall}
\longrightarrow
\text{singleton fibre / no intrinsic wall}.
\]

\begin{remark}[Intrinsic versus ambient criticality]
The three-anchor response image \(E_{\mathcal C}\) is a two-dimensional
embedded surface in \(\mathbb R^3\). The map cannot be a submersion onto the
ambient three-dimensional response space, yet it is a diffeomorphism onto its
intrinsic image. This illustrates why WR-III structural regularity is defined
relative to the realizable response image rather than by ambient surjectivity
alone.
\end{remark}

\begin{remark}[Prior-art status]
None of the trilateration facts above is claimed as new. Distance-based sensor
localization and unique localizability are classical; in planar network
localization, three or more noncollinear anchors play the standard role in
eliminating geometric ambiguity \cite{AspnesEtAl2006,AndersonEtAl2010}.
The WR-III contribution is the organization of this familiar inverse problem
into refinement, fibre-diameter, conditional-rank, and bifurcation-wall
diagnostics.
\end{remark}


\section{Prior-art boundary}

The audit places WR-III next to several mature mathematical lines.

\begin{itemize}
\item Set-valued analysis already studies upper/lower hemicontinuity and
inverse multifunctions. The open/closed-map criteria used here are
infrastructure rather than novelty claims.

\item Bifurcation and singularity theory use failure of local fibre
triviality as the definition of bifurcation or atypical values. For proper
submersions, Ehresmann gives local triviality; for nonproper maps, bifurcation
values at infinity may occur. In analytic/algebraic settings, critical loci of
augmented maps \((f,g)\), equivalently criticality of the added observable
along fibres of \(f\), are standardly organized by relative polar varieties.
Thus the conditional-rank set of Theorem~\ref{thm:conditional-rank} has direct
prior-art relatives and is not claimed as a new singularity-theory object.

\item Thom-type isotopy theorems and Hardt semialgebraic triviality provide
stratified or tame local-triviality frameworks beyond smooth manifolds.

\item Structural-identifiability geometry already uses parameter-fibre
dimension, degree, and irreducible components to diagnose identifiability.

\item Reeb-graph and Reeb-space constructions summarize changes in connected
components of fibres and their topology.
\end{itemize}

WR-III therefore does not claim invention of fibre geometry, bifurcation
sets, local triviality, or set-valued continuity. Its narrower target is the
WREH-specific organization of nonempty response fibres downstream of WR-I and
WR-II, together with protocol-refinement shrinkage, metric rigidity, and an
explicit claim firewall.

\section{Open obligations}

Before WR-III can become a preprint:

\begin{enumerate}
\item complete a source-level audit of the conditional-rank construction
against relative polar varieties and augmented-map singularity theory;
\item determine useful fibre invariants beyond diameter and prove which of
them remain monotone under protocol refinement;
\item formulate a quantitative distance-to-bifurcation notion only after a
response-space metric and regularity class are fixed;
\item extend the refinement theorem from smooth proper maps to a justified
Whitney/Thom or semialgebraic stratified setting;
\item extend the range-only localization benchmark to noisy ranges and a
bounded uncertainty model, and compare rigidity gain with standard localization
error bounds;
\item determine whether the rigidity-gain function
\(G_{\Bcal\mid\Acal}\) has stable or computable bounds in a useful model
class;
\item build an interactive demo showing the monotone rigidity/nonmonotone-wall
contrast.
\end{enumerate}

\section*{Conclusion}

WR-I identifies response-equivalence classes. WR-II determines whether
compatible responses admit global realizations. WR-III studies the geometry
remaining inside those nonempty inverse fibres.

The first exact chain is
\[
\text{refinement}
\Longrightarrow
\text{fibre inclusion}
\Longrightarrow
\text{nonincreasing diameter}.
\]
Compactness controls upper variation of fibres and fibre diameter. In the
smooth proper setting, regular values carry locally trivial fibres, leaving
critical values and attainable-response boundaries as natural containers for
model-relative realizability walls.

No wall introduced here is, without an additional physical bridge, a wall of
spacetime or of the Universe.

\begin{thebibliography}{99}

\bibitem{MalachevskyWRI}
A. A. Malachevsky,
\emph{Response-Defined World Spaces: Response Equivalence,
Finite-Resolution Separation, and Identifiability of Global Realizations},
WR-I, WREH (2026).
DOI: 10.5281/zenodo.23210258.

\bibitem{MalachevskyWRII}
A. A. Malachevsky,
\emph{Finite Consistency and Global World Realizability},
WR-II, WREH (2026).
DOI: 10.5281/zenodo.23224154.

\bibitem{AspnesEtAl2006}
J. Aspnes, T. Eren, D. K. Goldenberg, A. S. Morse, W. Whiteley,
Y. R. Yang, B. D. O. Anderson, and P. N. Belhumeur,
\emph{A Theory of Network Localization},
IEEE Transactions on Mobile Computing 5(12) (2006), 1663--1677.
DOI: 10.1109/TMC.2006.174.

\bibitem{AndersonEtAl2010}
B. D. O. Anderson, I. Shames, G. Mao, and B. Fidan,
\emph{Formal Theory of Noisy Sensor Network Localization},
SIAM Journal on Discrete Mathematics 24(2) (2010), 684--698.
DOI: 10.1137/100792366.

\bibitem{AubinFrankowska}
J.-P. Aubin and H. Frankowska,
\emph{Set-Valued Analysis},
Birkh\"auser, 1990.

\bibitem{Ehresmann}
C. Ehresmann,
\emph{Les connexions infinit\'esimales dans un espace fibr\'e
diff\'erentiable},
Colloque de Topologie, Bruxelles, 1950.

\end{thebibliography}

\end{document}
